\subsection{Binary Search Tree 기본 연산}
\begin{frame}{BST Property: Distinct-Key Convention}
모든 node $x$에 대해
\[
\forall k\in x.left:\ k<x.key,\qquad
\forall k\in x.right:\ k>x.key.
\]
\begin{alertblock}{Duplicate 정책}본문과 코드는 duplicate를 거부한다. count field 또는 항상 한쪽 삽입도 가능하지만 정책을 섞으면 안 된다.\end{alertblock}
\end{frame}
\begin{frame}{고정 BST 예제}
\centering\begin{tikzpicture}[level distance=9mm,level 1/.style={sibling distance=42mm},level 2/.style={sibling distance=22mm},level 3/.style={sibling distance=12mm}]
\node[st node]{15}
 child{node[st node]{6}child{node[st node]{3}child{node[st node]{2}}child{node[st node]{4}}}child{node[st node]{7}child[missing]{}child{node[st node]{13}child{node[st node]{9}}child{node[st node]{14}}}}}
 child{node[st node]{18}child{node[st node]{17}}child{node[st node]{20}}};
\end{tikzpicture}
\end{frame}
\begin{frame}{BST Inorder는 정렬 순서}
\[
2,3,4,6,7,9,13,14,15,17,18,20
\]
\begin{enumerate}\item left subtree key는 root보다 작다.\item root를 출력한다.\item right subtree key는 root보다 크다.\end{enumerate}
\stcheckpoint{모든 binary tree를 inorder하면 정렬되는가? \textbf{아니다.} BST invariant가 필요하다.}
\end{frame}
\begin{frame}{Recursive TreeSearch}
\small\begin{algorithmic}[1]
\Procedure{TreeSearch}{$x,k$}
\If{$x=\NIL$ \textbf{or} $k=x.key$}\State\Return $x$\EndIf
\If{$k<x.key$}\State\Return \Call{TreeSearch}{$x.left,k$}\EndIf
\State\Return \Call{TreeSearch}{$x.right,k$}
\EndProcedure
\end{algorithmic}
\end{frame}
\begin{frame}{Iterative TreeSearch}
\small\begin{algorithmic}[1]
\Procedure{IterativeTreeSearch}{$x,k$}
\While{$x\ne\NIL$ \textbf{and} $k\ne x.key$}
\If{$k<x.key$}\State $x\gets x.left$
\Else\State $x\gets x.right$\EndIf
\EndWhile
\State\Return $x$
\EndProcedure
\end{algorithmic}
\end{frame}
\begin{frame}{SEARCH 13: 실제 BST 경로}
\centering\begin{tikzpicture}[x=1.15cm,y=.62cm]
\node[font=\scriptsize,text=AlgoGray]at(-4.4,.2){full BST; inactive subtree는 탐색하지 않음};
\node[st inactive](n15)at(0,0){15};
\node[st inactive](n6)at(-3,-1){6}; \node[st inactive](n18)at(3,-1){18};
\node[st inactive](n3)at(-4.3,-2){3}; \node[st inactive](n7)at(-1.7,-2){7};
\node[st inactive](n17)at(2.3,-2){17}; \node[st inactive](n20)at(3.8,-2){20};
\node[st inactive](n2)at(-5,-3){2}; \node[st inactive](n4)at(-3.7,-3){4};
\node[st inactive](n13)at(-.5,-3){13};
\node[st inactive](n9)at(-1.2,-4){9}; \node[st inactive](n14)at(.2,-4){14};
\draw[st edge](n15)--(n6)(n15)--(n18)(n6)--(n3)(n6)--(n7)
 (n3)--(n2)(n3)--(n4)(n7)--(n13)(n13)--(n9)(n13)--(n14)
 (n18)--(n17)(n18)--(n20);
\only<2-|handout:1>{\draw[st active edge](n15)--(n6);}
\only<3-|handout:1>{\draw[st active edge](n6)--(n7);}
\only<4|handout:1>{\draw[st active edge](n7)--(n13);}
\only<1|handout:0>{\node[st current]at(n15){15};}
\only<2|handout:0>{\node[st done]at(n15){15};\node[st current]at(n6){6};}
\only<3|handout:0>{\node[st done]at(n15){15};\node[st done]at(n6){6};\node[st current]at(n7){7};}
\only<4|handout:1>{\node[st done]at(n15){15};\node[st done]at(n6){6};\node[st done]at(n7){7};\node[st current]at(n13){13};}
\only<1|handout:0>{\node[callout]at(0,-5){$13<15$: right subtree를 버리고 left로};}
\only<2|handout:0>{\node[callout]at(0,-5){$13>6$: left subtree를 버리고 right로};}
\only<3|handout:0>{\node[callout]at(0,-5){$13>7$: left/NIL을 버리고 right로};}
\only<4|handout:1>{\node[callout]at(0,-5){$13=13$: found};}
\end{tikzpicture}
\end{frame}
\begin{frame}{MINIMUM · MAXIMUM}
\begin{columns}[T]\column{.48\textwidth}
\begin{block}{Minimum}left child를 NIL까지 따라감: 예제에서 2\end{block}
\column{.48\textwidth}
\begin{block}{Maximum}right child를 NIL까지 따라감: 예제에서 20\end{block}
\end{columns}
\begin{itemize}\item time $\Theta(h)$\item empty tree의 Java API는 \texttt{OptionalInt.empty()} 반환\end{itemize}
\end{frame}
\begin{frame}{TreeMinimum Pseudocode}
\small\begin{algorithmic}[1]
\Procedure{TreeMinimum}{$x$}
\If{$x=\NIL$}\State\Return $\NIL$\EndIf
\While{$x.left\ne\NIL$}\State $x\gets x.left$\EndWhile
\State\Return $x$
\EndProcedure
\end{algorithmic}
\begin{block}{Maximum}left를 right로 바꾼 symmetric algorithm.\end{block}
\end{frame}
\begin{frame}{SUCCESSOR: 두 Case}
\begin{enumerate}
\item $x.right\ne\NIL$: $\operatorname{minimum}(x.right)$
\item right subtree 없음: parent를 올라가며 현재 node가 처음으로 ancestor의 left subtree에 놓이는 그 ancestor
\end{enumerate}
\begin{block}{정의}$x.key$보다 큰 key 중 가장 작은 key.\end{block}
\end{frame}
\begin{frame}{SUCCESSOR Animation}
\only<1|handout:0>{%
  \begin{columns}\column{.48\textwidth}\centering
  \begin{tikzpicture}[level distance=11mm,sibling distance=20mm]
  \node[st current]{15}child{node[st inactive]{6}}child{node[st done]{18}child{node[st node,draw=green!45!black,ultra thick,fill=green!12]{17}}child{node[st node]{20}}};
  \end{tikzpicture}
  \begin{block}{Case 1}$succ(15)=min(18\ subtree)=17$\end{block}
  \column{.48\textwidth}
  right subtree가 있으므로 그 안의 minimum
  \begin{align*}
  x.right &\ne \NIL,\\[-1mm]
  \operatorname{successor}(x)
    &= \operatorname{minimum}(x.right).
  \end{align*}
  \end{columns}}
\only<2|handout:1>{%
  \begin{columns}[T]
  \column{.58\textwidth}
  \begin{block}{Case 2: right subtree가 없을 때}
  원래 node가 ancestor의 left subtree에 속하는 첫 ancestor에 도달할 때까지
  upward 이동한다. 그 ancestor가 successor다.
  \end{block}
  \vspace{2mm}
  \centering
  \begin{tikzpicture}[level distance=12mm,sibling distance=20mm]
  \node[st node,draw=green!45!black,ultra thick,fill=green!12]{6}
    child{node[st node]{3}child[missing]{}child{node[st current]{4}}}
    child[missing]{};
  \end{tikzpicture}
  \column{.38\textwidth}
  \centering
  \begin{tikzpicture}[y=1.25cm]
  \node[st node,draw=green!45!black,ultra thick,fill=green!12](n6)at(0,0){6};
  \node[st node](n3)at(0,-1){3};
  \node[st current](n4)at(0,-2){4};
  \draw[st active edge](n4)--(n3);
  \draw[st active edge](n3)--(n6);
  \node[align=center,font=\small,text=green!45!black]at(0,-2.8)
    {$succ(4)=6$\\first qualifying ancestor};
  \end{tikzpicture}
  \end{columns}}
\end{frame}
\begin{frame}{SUCCESSOR 예제 검증}
\begin{columns}[T]
\column{.48\textwidth}
\centering
\begin{tikzpicture}[level distance=11mm,sibling distance=20mm]
\node[st current]{15}
  child{node[st inactive]{6}}
  child{node[st node]{18}
    child{node[st node,draw=green!45!black,ultra thick,fill=green!12]{17}}
    child{node[st node]{20}}};
\end{tikzpicture}
\begin{block}{Case 1}
$15$는 right subtree가 있다.\\
$succ(15)=\textcolor{green!45!black}{17}$
\end{block}
\column{.48\textwidth}
\centering
\begin{tikzpicture}[level distance=11mm,sibling distance=20mm]
\node[st node,draw=green!45!black,ultra thick,fill=green!12]{6}
  child{node[st node]{3}child[missing]{}child{node[st current]{4}}}
  child[missing]{};
\end{tikzpicture}
\begin{block}{Case 2}
$4$는 right subtree가 없다.\\
원래 node $4$를 left subtree에 포함하는 첫 ancestor는
\textcolor{green!45!black}{$6$}: $succ(4)=6$
\end{block}
\end{columns}
\vspace{1mm}
\centering\small
parent pointer가 없으면 root부터 search path를 기억하거나 다시 탐색한다.
\end{frame}
\begin{frame}{BST 기본 연산 Checkpoint}
\begin{enumerate}\item search(14)의 path는?\item minimum과 maximum은?\item predecessor(15)는?\item successor(20)는?\end{enumerate}
\begin{exampleblock}{답}$15\to6\to7\to13\to14$; 2와 20; 14; NIL.\end{exampleblock}
\end{frame}
